\documentclass{article}
\usepackage[utf8]{inputenc}
\usepackage[T1]{fontenc}
\usepackage{geometry}
\geometry{
    a4paper,
    left=2.5cm,
    right=2.5cm,
    top=2.5cm,
    bottom=2.5cm
}
\usepackage{graphicx}
\usepackage{float}
\usepackage{hyperref}
\usepackage[english]{babel}
\usepackage{amsmath}
\usepackage{amssymb}
\usepackage{array}
\usepackage{booktabs}
\usepackage{caption}
\usepackage{subcaption}
\usepackage{enumitem}
\usepackage{fancyhdr}
\usepackage{titlesec}
\usepackage{url}
\usepackage{textcomp}
\usepackage{xcolor}
\usepackage{listings}
\usepackage{multirow}
\usepackage{longtable}
\usepackage{tabularx}
\usepackage{times}

% Set up page style
\pagestyle{fancy}
\fancyhf{}
\rfoot{\thepage}
\renewcommand{\headrulewidth}{0pt}

% Title formatting
\titleformat{\chapter}[display]{\normalfont\bfseries}{\chaptertitle\thechapter}{0pt}{\Huge}
\titlespacing*{\chapter}{0pt}{50pt}{40pt}

% Document info
\title{Template Document}
\author{Academic Support Agent}
\date{\today}

\begin{document}
\maketitle

\title{Basic Styles for the NDSEG fellowship application}
\author{You}
\date{\today}
\maketitle

\section{Introduction}
This document serves as a README for understanding and using the provided LaTeX template (`ndseg.sty`). It demonstrates various nuances and configuration options available within the template.

\section{Adjustable Settings}
All settings described in this document can be adjusted directly in the included `ndseg.sty` file. This allows you to customize the template according to your specific needs.

\section{Font and Encoding}
The template uses the Times New Roman font and supports UTF-8 input encoding. These settings ensure that the document uses a standard, readable font and can handle a wide range of characters.

\section{Page Geometry}
The page margins are set to 0.7 inches:
\begin{verbatim}
\usepackage[margin=0.7in]{geometry}
\end{verbatim}

\section{Line Spacing}
Line spacing is set to single by default, providing a compact and readable layout:
\begin{verbatim}
\singlespacing
\end{verbatim}

\section{Section Titles}
Section titles are formatted to be inline, bold, and underlined. This ensures a clear and professional appearance:
\begin{verbatim}
\titleformat{\section}[runin]
  {\normalfont\normalsize\bfseries}
  {}{0em}{}
\titleformat{\subsection}[runin]
  {\normalfont\normalsize\bfseries}
  {}{0em}{}
\titleformat{\subsubsection}[runin]
  {\normalfont\normalsize\bfseries}
  {}{0em}{}
\end{verbatim}

\section{Spacing Around Section Titles}
The spacing before and after section headings is adjusted to provide a clean look:
\begin{verbatim}
\titlespacing{\section}{0pt}{*0}{0.7em}
\titlespacing{\subsection}{0pt}{*0}{0.5em}
\titlespacing{\subsubsection}{0pt}{*0}{0.3em}
\end{verbatim}

\section{Paragraph Formatting}
Paragraph indentation is removed to create a modern look, and word spacing is adjusted for better readability:
\begin{verbatim}
\setlength{\parindent}{0pt}
\spaceskip=0.3em
\xspaceskip=0.3em
\end{verbatim}

\section{Figures}
The template includes several customizations for figures:
\begin{itemize}
  \item The default figure directories are set to `./figures/` and `./images/`:
    \begin{verbatim}
    \graphicspath{{./figures/}{./images/}}
    \end{verbatim}
  \item A figure default figure width is defined as 0.5 of the text width, which can be easily changed:
    \begin{verbatim}
% Set default width for graphics
% change me to change the default width
\graphicxGlobalKeys{width=0.5\textwidth}
    \end{verbatim}
\end{itemize}
This default can also be overwritten in the typical way by specifying
A different value for the width or scale parameters.
Here is an example of including a PNG image with the default width:

\begin{figure}[h!]
  \centering
  \includegraphics{placeholder.jpg} % Uses default width
  \caption{An example PNG image with default width.}
  \label{fig:example-png}
\end{figure}

Here is an example of including a JPG image with a custom width:

\begin{figure}[h!]
  \centering
  \includegraphics[width=2em]{placeholder.jpg} % Overrides default width
  \caption{An example JPG image with custom width.}
  \label{fig:example-jpg}
\end{figure}

Text can also wrap around a figure by using the \begin{verbatim}wrapfigure\end{verbatim} environment as shown in \figref{fig:example-wrapping}.

\begin{wrapfigure}{r}{0.5\textwidth}
  \begin{center}
    \includegraphics[width=0.48\textwidth]{placeholder.png}
  \end{center}
  \caption{Text wraps around}
  \label{fig:example-wrapping}
\end{wrapfigure}

\section{Hyphenation and Line Breaks}
To improve text justification and avoid overfull lines, hyphenation penalties and tolerances are adjusted:
\begin{verbatim}
\hyphenpenalty=5000
\tolerance=1000
\emergencystretch=.4em
\end{verbatim}

\section{Page Style}
The document uses a plain page style to remove the title page, header, and footer:
\begin{verbatim}
\pagestyle{plain}
\end{verbatim}

\section{Bibliography}
The bibliography style is set to `plain`, but you can change it to other styles like `abbrv` or `alpha` as needed:
\begin{verbatim}

\end{verbatim}

The bibligoraphy can be included with the bibliography command.
\begin{verbatim}
    \bibliography
\end{verbatim}

\section{Custom Commands}
The template includes custom commands for TODOs and comments, which are ignored during compilation:
\begin{verbatim}
\TODO{Fix This} % Inline TODO comments
\TODO{This is ignored} % Inline comments
\end{verbatim}

As well as a few quality of life commands such as
\begin{verbatim}
    \figref{fig:example} % Append Fig. before the reference number
\end{verbatim}

\section{Draft Mode}
For draft versions, line numbers can be enabled by setting the `draft` boolean to true:
\begin{verbatim}
\newboolean{draft}
\ifthenelse{\boolean{draft}}{
  \linenumbers
}{}
\end{verbatim}



\section{Conclusion}
All these settings can be adjusted directly in the included `ndseg.sty` file. This allows you to customize the template according to your specific needs. Simply open the `ndseg.sty` file and modify the relevant sections to suit your preferences.
\end{document}